\documentclass[11pt,a4paper]{article}
\usepackage[margin=1in]{geometry}
\usepackage{graphicx}
\usepackage{booktabs}
\usepackage{array}
\usepackage{longtable}
\usepackage{enumitem}
\usepackage{hyperref}
\usepackage{xcolor}
\usepackage{fancyhdr}
\usepackage{listings}
\usepackage{float}

% Define colors for warnings and notes
\definecolor{warningcolor}{RGB}{204,0,0}
\definecolor{notecolor}{RGB}{0,102,204}
\definecolor{tipcolor}{RGB}{0,153,0}

% Setup code listing style
\lstset{
    basicstyle=\ttfamily\small,
    frame=single,
    breaklines=true
}

% Custom commands for special formatting
\newcommand{\warning}[1]{\textcolor{warningcolor}{\textbf{WARNING:} #1}}
\newcommand{\note}[1]{\textcolor{notecolor}{\textbf{Note:} #1}}
\newcommand{\tip}[1]{\textcolor{tipcolor}{\textbf{Tip:} #1}}

\pagestyle{fancy}
\fancyhf{}
\rhead{BBSHD to Baserunner Z9 Wiring Guide}
\lhead{\leftmark}
\rfoot{Page \thepage}

\title{\textbf{Rewiring a Grin Baserunner V6\_Z9 to a Bafang BBSHD Motor}\\
\large Complete Installation and Configuration Guide}
\author{Technical Reference Document}
\date{\today}

\begin{document}

\maketitle
\tableofcontents
\newpage

\section{Important Safety Notice}

\warning{This modification involves cutting factory connectors and opening the motor/controller -- it voids all warranties. Always disconnect the battery and capacitors before working on wiring, and double-check all connections for correct polarity and insulation before applying power. Proceed only if you are comfortable with crimping and ebike electronics.}

\section{Overview}

The goal is to replace the BBSHD's internal controller with a Grin Baserunner V6 (Z9 model) external controller, creating a clean color-matched harness between the Baserunner and the BBSHD motor. The modification involves:

\begin{enumerate}
    \item Removing the stock BBSHD controller and its mating connector, exposing the motor's phase wires and sensor wires
    \item Cutting off the Baserunner's Higo Z910 plug (the 9-pin motor connector) and terminating its wires with new connectors
    \item Attaching Anderson Powerpole PP45 connectors to the three phase wires (Yellow, Green, Blue) for high-current connections
    \item Attaching a 6-pin JST-SM connector to the small hall/temperature sensor wires (5V, GND, Hall A, Hall B, Hall C, Thermistor)
\end{enumerate}

This will allow the Baserunner to plug directly into the BBSHD motor leads via robust, serviceable connectors. The result is a harness with color-to-color wiring between controller and motor for clarity.

\section{Wiring Pinout Mapping}

Before cutting any wires, it is critical to understand the wire color mapping between the Baserunner V6\_Z9 motor cable and the BBSHD motor's leads. The Baserunner's ``Z9'' motor cable contains 3 large phase conductors and 6 smaller signal wires. The BBSHD's motor has a matching set of wires once the internal controller is removed.

\subsection{Connection Mapping Table}

\begin{table}[H]
\centering
\caption{Baserunner to BBSHD Wiring Connections}
\label{tab:wiring}
\begin{tabular}{p{2.5cm}p{3.5cm}p{3.5cm}p{3.5cm}}
\toprule
\textbf{Signal} & \textbf{Baserunner V6\_Z9 Wire} & \textbf{BBSHD Motor Wire} & \textbf{Connector Type} \\
\midrule
Phase U & Yellow (large) & Yellow & Anderson PP45 (Yellow) \\
Phase V & Green (large) & Green & Anderson PP45 (Green) \\
Phase W & Blue (large) & Blue & Anderson PP45 (Blue) \\
Hall Sensor A & Yellow (small) & Gray & JST-SM 6-pin \\
Hall Sensor B & Green (small) & White & JST-SM 6-pin \\
Hall Sensor C & Blue (small) & Blue & JST-SM 6-pin \\
Hall/Logic +5V & Red (small) & Red & JST-SM 6-pin \\
Hall/Logic GND & Black (small) & Black & JST-SM 6-pin \\
Motor Temp & White (small) & Thermistor lead & JST-SM 6-pin \\
\bottomrule
\end{tabular}
\end{table}

\subsection{Important Wiring Notes}

\begin{itemize}
    \item The three phase wires match colors one-to-one (yellow--yellow, green--green, blue--blue)
    \item The BBSHD's hall sensor outputs use Blue, White, and Gray wires (rather than yellow/green/blue) -- these correspond to the Baserunner's hall inputs Blue, Green, and Yellow respectively
    \item The BBSHD's +5V (red) and ground (black) wires for the halls line up with the Baserunner's red and black
    \item For motor temperature, the Baserunner's white signal wire is intended for a thermistor input
    \item Bafang BBSHD motors typically include a 10k$\Omega$ NTC thermistor embedded in the windings
\end{itemize}

\section{Connectors and Wiring Harness}

\subsection{Anderson Powerpole PP45}

\begin{itemize}
    \item \textbf{Purpose:} High-current phase connections
    \item \textbf{Rating:} Up to 45A continuous
    \item \textbf{Features:} Genderless, modular connector
    \item \textbf{Color coding:} Yellow, green, and blue housings to match phase wire colors
    \item \textbf{Assembly:} Can be dovetailed together into a 3-pin block
\end{itemize}

\subsection{JST-SM 6-pin}

\begin{itemize}
    \item \textbf{Purpose:} Hall sensors and temperature sensor connections
    \item \textbf{Configuration:} 6 positions for 5 hall/therm signals plus ground and 5V
    \item \textbf{Convention:} Controller side uses male JST housing (with female pins), motor side uses female JST housing (with male pins)
    \item \textbf{Feature:} Female JST has a latch for secure connection
\end{itemize}

\section{Step-by-Step Wiring Assembly}

\subsection{Step 1: Remove Original Components}

\subsubsection{Detach the BBSHD Controller}
\begin{enumerate}
    \item Remove the stock Bafang BBSHD internal controller from the motor housing
    \item Unscrew and gently unplug its connectors
    \item Expose the internal wiring harness:
    \begin{itemize}
        \item Three thick phase wires (Yellow, Green, Blue)
        \item Cluster of thinner wires for hall sensors (Red, Black, Blue, White, Gray)
        \item Possible separate pair for temperature or PAS sensor
    \end{itemize}
    \item Disconnect or desolder the old controller from these wires as needed
\end{enumerate}

\subsubsection{Cut the Baserunner Z910 Plug}
\begin{enumerate}
    \item Ensure the controller is completely powered off
    \item Cut off the Z910 plug, leaving maximum cable length attached to the controller
    \item Strip back outer cable sheath to expose 9 individual wires:
    \begin{itemize}
        \item 3 larger gauge wires (Y, G, B for phases)
        \item 6 smaller wires (Red, Black, Yellow, Green, Blue, White)
    \end{itemize}
    \item Strip 3--4 mm of insulation from each wire end
    \item \tip{Do not tin wires that will be crimped as it can lead to a loose connection}
\end{enumerate}

\subsection{Step 2: Terminate Phase Wires with Anderson PP45}

\subsubsection{Prepare the Phase Leads}
\begin{enumerate}
    \item Identify the six phase wire ends (three from BBSHD, three from Baserunner)
    \item Strip approximately 5 mm of insulation from each wire
    \item Ensure copper strands are neatly twisted (typical gauge: 12--14 AWG)
\end{enumerate}

\subsubsection{Crimp the Contacts}
\begin{enumerate}
    \item Insert each stripped phase wire into a PP45 contact pin
    \item Use proper Anderson crimp tool to crimp the barrel tightly
    \item Verify secure connection with a light tug test
    \item Repeat for all six phase wire ends
\end{enumerate}

\subsubsection{Insert into Housings}
\begin{enumerate}
    \item Insert crimped contacts into color-matched PP45 housings:
    \begin{itemize}
        \item Yellow wire $\rightarrow$ Yellow housing
        \item Green wire $\rightarrow$ Green housing
        \item Blue wire $\rightarrow$ Blue housing
    \end{itemize}
    \item Ensure contact locks with audible click
    \item \note{Maintain consistent contact orientation (open side facing up toward ``A'' marking)}
\end{enumerate}

\subsubsection{Form 3-Pin Blocks}
\begin{enumerate}
    \item Join individual PP45 housings using dovetail slide rails
    \item Create two 3-position blocks (motor side and controller side)
    \item Verify color sequence matches on both blocks
    \item Test mating of blocks to confirm proper alignment
\end{enumerate}

\subsection{Step 3: Terminate Hall/Temp Wires with JST-SM 6-Pin}

\subsubsection{Identify Motor Sensor Wires}
\begin{enumerate}
    \item Locate on BBSHD motor side:
    \begin{itemize}
        \item Red (+5V)
        \item Black (Ground)
        \item Three hall signal wires (Blue, White, Gray)
        \item Thermistor lead (if present)
    \end{itemize}
    \item Plan for 6 total connections
    \item Isolate and secure any unused PAS sensor wires
\end{enumerate}

\subsubsection{Strip and Prepare Wires}
\begin{enumerate}
    \item Strip 2--3 mm of insulation from each small wire
    \item Handle carefully (typical gauge: 24--26 AWG)
    \item Twist strands to ensure neat insertion into pins
\end{enumerate}

\subsubsection{Crimp JST Pins}
\begin{enumerate}
    \item Use precision crimp tool for JST-SM pins
    \item Position wire with:
    \begin{itemize}
        \item Stripped conductor under smaller crimp wings
        \item Insulation under larger wings for strain relief
    \end{itemize}
    \item Crimp until tool fully closes, creating ``B''-shaped wrap
    \item Start with +5V (red) and GND (black) for easier tracking
\end{enumerate}

\subsubsection{Insert Wires into JST Housing}

\begin{table}[H]
\centering
\caption{JST Pin Assignment (Example Configuration)}
\begin{tabular}{cll}
\toprule
\textbf{Pin} & \textbf{Controller Side} & \textbf{Motor Side} \\
\midrule
1 & Red (+5V) & Red (+5V) \\
2 & Yellow (Hall A) & Gray (Hall A) \\
3 & Green (Hall B) & White (Hall B) \\
4 & Blue (Hall C) & Blue (Hall C) \\
5 & White (Temp) & Thermistor \\
6 & Black (GND) & Black (GND) \\
\bottomrule
\end{tabular}
\end{table}

\begin{enumerate}
    \item Motor side: Use female JST housing with male pins
    \item Controller side: Use male JST housing with female pins
    \item Insert pins until they click and lock
    \item Verify proper order for color-to-color mapping
\end{enumerate}

\subsection{Step 4: Insulation and Strain Relief}

\subsubsection{Insulate Exposed Conductors}
\begin{itemize}
    \item Verify Anderson housings fully cover phase contacts
    \item Add heat-shrink tubing over wire-to-housing junctions if desired
    \item Ensure no bare metal is exposed on JST connections
\end{itemize}

\subsubsection{Re-install Grommets/Seals}
\begin{itemize}
    \item Re-use or replace rubber grommet at motor wire exit
    \item Feed new phase and hall wires through properly
    \item Seat grommet securely to prevent water ingress
\end{itemize}

\subsubsection{Internal Strain Relief}
\begin{itemize}
    \item Secure wires inside motor with zip ties or original clamps
    \item Keep wires away from moving parts (rotor, gears)
    \item Bundle hall wires neatly
    \item Consider wrapping JST connector in foam if housed internally
\end{itemize}

\subsubsection{Connector Positioning}
\begin{itemize}
    \item Position Anderson and JST connectors for accessibility
    \item Allow sufficient slack for steering/suspension movement
    \item Consider using split loom or braided sleeve for protection
    \item Ensure connectors are not under tension
\end{itemize}

\section{Post-Install Configuration Checklist}

\subsection{Hardware Connections}

\subsubsection{Connect Main-9 SuperHarness}
\begin{enumerate}
    \item Plug SuperHarness 9-pin connector into Baserunner's Main9 cable
    \item Secure with threaded connection
    \item Connect SW102 display to SuperHarness Bafang display connector (green 5-pin Higo)
    \item Connect additional peripherals:
    \begin{itemize}
        \item E-brake cutoff sensors
        \item Throttle
        \item PAS sensor (if using external)
    \end{itemize}
\end{enumerate}

\subsection{Software Configuration}

\subsubsection{Load Default Profile}
\begin{enumerate}
    \item Launch Grin's Phaserunner/Baserunner Suite on PC
    \item Connect to Baserunner via USB (TTL-USB programming cable)
    \item Load profile: ``V6\_Baserunner\_w\_SuperH'' or similar
    \item This enables:
    \begin{itemize}
        \item Digital third-party display support
        \item TX/RX communication to SW102
        \item Proper throttle/brake input mapping
    \end{itemize}
\end{enumerate}

\subsubsection{Verify Display Communication}
\begin{enumerate}
    \item Power on system (battery connected, Baserunner keyed on)
    \item Confirm SW102 display powers up
    \item Test assist level changes via display buttons
    \item Set initial assist to 0 or low level for safety
\end{enumerate}

\subsection{Motor Calibration}

\subsubsection{Motor Hall Detection \& Autotune}
\begin{enumerate}
    \item Secure bike (wheel off ground or chain removed)
    \item In software Motor tab, select ``Hall Sensor motor''
    \item Click ``Start Autotune'' or ``Detect Motor Params''
    \item Motor will emit pulses and rotate slowly
    \item Process measures:
    \begin{itemize}
        \item Hall timing
        \item Resistance (Rs)
        \item Inductance (L)
        \item EMF constant (Kv)
    \end{itemize}
\end{enumerate}

\subsubsection{Expected Autotune Results}
\begin{itemize}
    \item Effective Pole Pairs: $\sim$88 (due to internal gearing)
    \item No-load speed constant: $\sim$3.3 RPM/V
    \item Motor should run smoothly after calibration
\end{itemize}

\subsection{Thermal Protection Setup}

\subsubsection{Configure Thermal Rollback}
\begin{enumerate}
    \item Navigate to Temperature (Thermistor) settings tab
    \item Select: ``NTC 10k Thermistor (Beta 3950)''
    \item Set temperature thresholds:
    \begin{itemize}
        \item Start reduction: 100°C
        \item Significant cutback: 120°C
    \end{itemize}
    \item Enable rollback feature
\end{enumerate}

\subsubsection{Verify Temperature Reading}
\begin{itemize}
    \item Check ambient temperature reading (should be 20--25°C when cool)
    \item If reading is incorrect:
    \begin{itemize}
        \item Newer BBSHD units (2020+) may use different thermistor
        \item Consider replacing with standard 10k$\Omega$ NTC if needed
    \end{itemize}
\end{itemize}

\subsection{Final Parameter Adjustments}

\subsubsection{Current Limits}
\begin{itemize}
    \item \textbf{Battery current:} 25--30A (based on battery capability)
    \item \textbf{Phase current:} 60--80A (for strong performance)
    \item \warning{Higher currents generate more heat}
\end{itemize}

\subsubsection{Control Settings}
\begin{itemize}
    \item \textbf{Throttle:} Configure input mapping and mode (torque/speed)
    \item \textbf{E-brakes:} Set regen level or cutoff behavior
    \item \textbf{PAS:} Configure if using external sensor (pulses, assist modes)
    \item \textbf{Speed source:} Connect external wheel speed sensor to SW102 or CA3
\end{itemize}

\subsection{Testing and Validation}

\subsubsection{Final System Checks}
\begin{enumerate}
    \item Review diagnostics in software with system powered
    \item Rotate cranks by hand -- verify hall sensors toggle
    \item Test throttle response (wheel free)
    \item Confirm brake cutoffs function
\end{enumerate}

\subsubsection{Test Ride Protocol}
\begin{enumerate}
    \item Reassemble BBSHD motor cover
    \item Secure all cables
    \item Start with low assist level or minimal throttle
    \item Monitor for:
    \begin{itemize}
        \item Smooth, quiet operation (FOC controller benefit)
        \item Temperature readings during hill climbs
        \item Thermal rollback engagement if needed
    \end{itemize}
    \item Check connectors for heat or loosening after first rides
\end{enumerate}

\section{Troubleshooting Guide}

\subsection{Common Issues and Solutions}

\begin{table}[H]
\centering
\caption{Troubleshooting Reference}
\begin{tabular}{p{4cm}p{9cm}}
\toprule
\textbf{Issue} & \textbf{Solution} \\
\midrule
Motor stutters or fails to sync & Verify hall sensor wire mapping (Gray$\rightarrow$Yellow, White$\rightarrow$Green, Blue$\rightarrow$Blue) \\
\midrule
Display doesn't power on & Check Baserunner's orange ``Key'' wire connection to battery positive \\
\midrule
Incorrect temperature reading & Verify thermistor type (10k$\Omega$ NTC), check wiring, consider replacement if needed \\
\midrule
No speed reading & Connect external wheel speed sensor to SW102 or CA3 \\
\midrule
Motor runs rough & Re-run autotune procedure, verify phase connections \\
\bottomrule
\end{tabular}
\end{table}

\section{Maintenance Notes}

\begin{itemize}
    \item Anderson and JST connections allow easy disconnection for service
    \item Monitor connectors during first 100km for any issues
    \item Keep JST connections protected from moisture
    \item Periodically check temperature readings for system health
    \item Save configuration files after any adjustments
\end{itemize}

\section{References}

\begin{enumerate}
    \item Grin Technologies, \textit{Baserunner V6 User Manual}, motor connector pinouts
    \item Grin Technologies, \textit{Connectors Overview}, Anderson Powerpoles and JST-SM for ebike wiring
    \item Endless-Sphere Forum, BBSHD hall/phase color mapping discussions
    \item Endless-Sphere Forum, BBSHD temperature sensor specifications (10k thermistor vs newer sensors)
    \item Reddit r/ebikes, BBSHD external controller experiences and profiles
\end{enumerate}

\vspace{1cm}
\begin{center}
\rule{0.5\textwidth}{0.4pt}\\[0.5cm]
\textit{End of Document}
\end{center}

\end{document}